% ══════════════════════════════════════════════════════════════════════════════
% DMAPO: Data-centric Multi-Agent Preference Optimization
% Experiment Results & Methods
% ══════════════════════════════════════════════════════════════════════════════

\documentclass[10pt,twocolumn]{article}

% ── Packages ──────────────────────────────────────────────────────────────────
\usepackage[utf8]{inputenc}
\usepackage[T1]{fontenc}
\usepackage{amsmath,amssymb,amsfonts}
\usepackage{booktabs}
\usepackage{multirow}
\usepackage{graphicx}
\usepackage{xcolor}
\usepackage{algorithm}
\usepackage{algpseudocode}
\usepackage{hyperref}
\usepackage{enumitem}
\usepackage{caption}
\usepackage{subcaption}
\usepackage[margin=1in]{geometry}
\usepackage{times}
\usepackage{microtype}

% ── Custom commands ───────────────────────────────────────────────────────────
\newcommand{\ours}{\textsc{DMAPO}}
\newcommand{\best}[1]{\textbf{#1}}
\newcommand{\second}[1]{\underline{#1}}
\newcommand{\up}{$\uparrow$}
\newcommand{\down}{$\downarrow$}
\newcommand{\cmark}{\ding{51}}
\newcommand{\xmark}{\ding{55}}

\begin{document}

% ══════════════════════════════════════════════════════════════════════════════
%
%   PART I — METHOD
%
% ══════════════════════════════════════════════════════════════════════════════

\section{Method}
\label{sec:method}

We present \ours{} (\textbf{D}ata-centric \textbf{M}ulti-\textbf{A}gent \textbf{P}reference \textbf{O}ptimization), a pipeline that constructs high-quality preference data through multi-agent evaluation and strict confidence gating, then trains a policy model on the curated signal.
Unlike prior work that relies on large, noisy preference corpora, \ours{} takes a \emph{data-centric} approach: we invest computation in data quality rather than data quantity, retaining only examples on which multiple independent judges reach strong consensus.
The method consists of six stages (\S\ref{sec:prompt}--\S\ref{sec:training}), illustrated in Figure~\ref{fig:pipeline} and summarized in Algorithm~\ref{alg:pipeline}.


\subsection{Overview}

The core intuition behind \ours{} is that preference optimization is fundamentally limited by the quality of the training signal.
Standard preference datasets~\cite{cui2023ultrafeedback,bai2022training} contain substantial annotation noise: ambiguous preference pairs, mislabeled examples, and borderline cases where even human annotators would disagree.
Training on such data forces the model to fit contradictory gradients, degrading alignment quality.

\ours{} addresses this by:
\begin{enumerate}[leftmargin=*,itemsep=1pt,topsep=2pt]
  \item \textbf{Generating diverse candidates} from the policy backbone itself, ensuring the training distribution matches the model's own generation capacity;
  \item \textbf{Scoring each candidate independently} along three orthogonal quality dimensions using specialized LLM judge agents;
  \item \textbf{Detecting and penalizing reasoning flaws} via a dedicated process critic that identifies circular, unsupported, or contradictory logic;
  \item \textbf{Aggressively filtering} via a confidence gate that retains only examples with high inter-judge agreement and extreme quality scores;
  \item \textbf{Training on the curated signal} using Kahneman-Tversky Optimization (KTO)~\cite{ethayarajh2024kto}, which naturally accommodates the binary desirable/undesirable labels produced by our gating mechanism.
\end{enumerate}


\subsection{Stage 1: Prompt Collection and Normalization}
\label{sec:prompt}

We source instruction prompts from two complementary datasets:
\textbf{UltraFeedback}~\cite{cui2023ultrafeedback} (10{,}000 prompts covering diverse instruction-following tasks) and \textbf{HelpSteer2}~\cite{wang2024helpsteer2} (5{,}000 prompts emphasizing helpfulness and steerability).
We extract the raw instruction field from each dataset, deduplicate by exact string match, and normalize formatting to produce a unified prompt pool.
The combined pool contains 14{,}272 unique prompts, which we split into training (13{,}559; 95\%) and validation (713; 5\%) subsets using a fixed random seed for reproducibility.

The diversity of the source pool is intentional: by combining broadly-scoped UltraFeedback instructions with the steered, quality-focused HelpSteer2 prompts, we expose the downstream judges to a wide range of instruction types, reducing domain-specific bias in the scoring stage.


\subsection{Stage 2: Candidate Response Generation}
\label{sec:generation}

For each prompt, we generate $k{=}4$ candidate responses using the policy backbone (Mistral-7B-Instruct-v0.2) with nucleus sampling ($T{=}0.8$, $p{=}0.95$).
Crucially, we generate from the \emph{same model} that will later be fine-tuned.
This on-policy generation strategy ensures that the preference signal is calibrated to the model's own output distribution, avoiding the distribution mismatch that arises when training on preferences over outputs from a different, typically stronger model~\cite{gao2023scaling}.

To maximize throughput, we tile each prompt $k$ times and process all candidates in a single batched \texttt{model.generate()} call per prompt group.
Combined with \texttt{torch.compile} for compute-graph optimization, this yields a generation rate of approximately 2{,}000 candidates per GPU-hour on NVIDIA RTX PRO 6000 Blackwell GPUs.
The full candidate pool comprises 54{,}236 training and 2{,}852 validation candidates.


\subsection{Stage 3: Multi-Agent Scoring}
\label{sec:scoring}

Each candidate response is evaluated independently by a panel of three specialized LLM judge agents, each instantiated from Qwen3-8B~\cite{qwen2025qwen3} in deterministic mode (\texttt{temperature}$=$1.0, \texttt{do\_sample}$=$False):

\paragraph{Helpfulness Judge.}
Evaluates whether the response directly and completely addresses the user's instruction, assigning a score from 1 (not helpful) to 10 (extremely helpful).

\paragraph{Factuality Judge.}
Assesses the factual accuracy of the response content, scoring from 1 (many factual errors) to 10 (fully accurate).

\paragraph{Conciseness Judge.}
Evaluates response efficiency---whether the response conveys the necessary information without unnecessary verbosity or padding, scoring from 1 (very verbose) to 10 (perfectly concise).

Each judge operates through a structured chat prompt with a dimension-specific system instruction (see Appendix for exact prompts) and outputs a single integer score followed by an optional brief justification.
The scores are parsed via regex to extract the numeric value, with a fallback to the midpoint score of 5.5 when parsing fails.

The use of orthogonal evaluation dimensions---rather than a single holistic quality score---serves two purposes.
First, it decomposes quality assessment into tractable sub-problems, reducing the cognitive burden on each judge and improving scoring reliability.
Second, inter-judge variance becomes a meaningful signal for downstream confidence gating: high variance indicates that the candidate performs well on some dimensions but poorly on others, suggesting ambiguous quality that should be excluded from training.


\subsection{Stage 4: Process Critic}
\label{sec:critic}

Beyond surface-level quality, we deploy a \textbf{process critic}---also instantiated from Qwen3-8B---that specifically targets reasoning quality.
The critic reads the instruction-response pair and determines whether the response contains \emph{flawed}, \emph{circular}, or \emph{unsupported} reasoning.
It outputs a binary verdict (\texttt{YES}/\texttt{NO}), a failure type classification (circular $|$ unsupported $|$ contradictory), and a brief explanation.

When the critic detects flawed reasoning, a penalty is applied to the candidate's aggregate score:
\begin{equation}
  s_{\text{final}} = \bar{s}_{\text{judges}} - \alpha \cdot \mathbb{1}[\text{flawed}] \cdot \bar{s}_{\text{judges}}
  \label{eq:critic_penalty}
\end{equation}
where $\bar{s}_{\text{judges}} = \frac{1}{J}\sum_{j=1}^{J} s_j$ is the mean score across $J{=}3$ judges, $\mathbb{1}[\text{flawed}]$ is the critic's binary indicator, and $\alpha{=}0.15$ is the penalty weight.
This multiplicative penalty is intentionally moderate: it reduces the final score enough to push borderline-desirable candidates below the gating threshold without completely overriding the judge consensus.


\subsection{Stage 5: Confidence Gating}
\label{sec:gating}

The confidence gate is the key mechanism that distinguishes \ours{} from standard data collection pipelines.
Given the multi-agent scores and critic-adjusted final score for each candidate, we apply a two-stage filter:

\paragraph{Variance Gate.}
We compute the inter-judge score variance $\sigma^2 = \text{Var}(s_1, s_2, s_3)$ across the three judge dimensions.
Candidates with $\sigma^2 > \tau_{\text{var}}$ (where $\tau_{\text{var}}{=}2.5$) are discarded as ambiguous---the judges fundamentally disagree on the candidate's quality, so any label assignment would be unreliable.

\paragraph{Quality Gate.}
For candidates passing the variance gate, we assign binary labels based on extreme-score thresholds:
\begin{equation}
  \text{label} =
  \begin{cases}
    \texttt{desirable}   & \text{if } s_{\text{final}} \geq \tau_{\text{des}} \\
    \texttt{undesirable} & \text{if } s_{\text{final}} \leq \tau_{\text{und}} \\
    \texttt{discard}     & \text{otherwise}
  \end{cases}
  \label{eq:gating}
\end{equation}
where $\tau_{\text{des}}{=}7.0$ and $\tau_{\text{und}}{=}4.0$.
Candidates falling in the ``gray zone'' between thresholds are discarded---they carry weak or ambiguous preference signal.

This aggressive filtering yields a \textbf{3.45\% acceptance rate} on the training set, retaining only 1{,}871 examples from 54{,}236 scored candidates.
The resulting dataset is highly polarized: desirable examples have a mean final score of $9.23 \pm 1.09$ (range 7--10), while undesirable examples score $2.42 \pm 1.10$ (range 1--4), producing a large quality gap of approximately 6.8 points.
This polarization ensures that every training example carries a strong, unambiguous preference signal.


\subsection{Stage 6: Policy Training}
\label{sec:training}

We train the policy model using \textbf{Kahneman-Tversky Optimization (KTO)}~\cite{ethayarajh2024kto} rather than the more common DPO~\cite{rafailov2023dpo}.
The choice of KTO is motivated by two factors:
\begin{enumerate}[leftmargin=*,itemsep=1pt,topsep=2pt]
  \item \textbf{Natural fit with binary labels.} Our gating mechanism produces binary (desirable/undesirable) labels for individual responses, which directly maps to KTO's training format. DPO requires constructing explicit (chosen, rejected) pairs from the same prompt, which discards candidates and introduces artificial pairing artifacts.
  \item \textbf{Loss-aversion asymmetry.} KTO incorporates the insight from prospect theory~\cite{kahneman1979prospect} that losses loom larger than gains, applying asymmetric weighting to desirable and undesirable examples. This aligns well with our quality-gated data, where both the ``clearly good'' and ``clearly bad'' signals are independently informative.
\end{enumerate}

\paragraph{Model and Adapter Configuration.}
We fine-tune Mistral-7B-Instruct-v0.2~\cite{jiang2023mistral} (7.24B parameters) using Low-Rank Adaptation (LoRA)~\cite{hu2022lora} with rank $r{=}16$, scaling factor $\alpha{=}32$, and dropout of 0.05.
LoRA adapters are applied to all attention projections (Q, K, V, O) and feedforward projections (gate, up, down), totaling approximately 160M trainable parameters ($\sim$2.2\% of the full model).

\paragraph{Training Hyperparameters.}
We use AdamW~\cite{loshchilov2019decoupled} with fused kernels, a cosine learning rate schedule with 5\% warmup, peak learning rate of $5 \times 10^{-5}$, and an effective batch size of 16 (2 per-device $\times$ 8 gradient accumulation steps).
Training runs for 1 epoch in bfloat16 precision.
The KTO-specific preference strength parameter is $\beta{=}0.1$.

\paragraph{Training Dynamics.}
Despite training on only 1{,}871 examples (5$\times$ fewer than the baselines), the model converges to a final loss of 0.037 with a reward margin of 10.96.
The high reward margin---compared to 1.88 for a DPO baseline trained on 10{,}000 UltraFeedback pairs---reflects the extremely clean training signal: with no contradictory labels to fit, the model can decisively separate desirable from undesirable outputs.


\begin{algorithm}[t]
\caption{\ours{} Pipeline}
\label{alg:pipeline}
\begin{algorithmic}[1]
\Require Prompt pool $\mathcal{P}$, policy backbone $\pi_\theta$, judges $\{J_1, \ldots, J_M\}$, critic $C$, thresholds $\tau_\text{var}, \tau_\text{des}, \tau_\text{und}$
\Ensure Aligned policy $\pi_{\theta^*}$
\State \textbf{Stage 1:} Collect and normalize prompts $\mathcal{P} \leftarrow \mathcal{P}_\text{UF} \cup \mathcal{P}_\text{HS2}$
\For{each prompt $x \in \mathcal{P}$} \Comment{Stage 2}
  \State Generate $k$ candidates: $\{y_1, \ldots, y_k\} \sim \pi_\theta(\cdot \mid x)$
\EndFor
\For{each candidate $(x, y_i)$} \Comment{Stage 3}
  \State Score: $s_j \leftarrow J_j(x, y_i)$ for $j = 1, \ldots, M$
  \State Compute mean: $\bar{s} \leftarrow \frac{1}{M}\sum_j s_j$
  \State Compute variance: $\sigma^2 \leftarrow \text{Var}(s_1, \ldots, s_M)$
\EndFor
\For{each candidate $(x, y_i)$} \Comment{Stage 4}
  \State Critic: $p \leftarrow C(x, y_i)$ \Comment{$p \in \{0, 1\}$}
  \State $s_\text{final} \leftarrow \bar{s} - \alpha \cdot p \cdot \bar{s}$
\EndFor
\For{each candidate $(x, y_i)$} \Comment{Stage 5: Gating}
  \If{$\sigma^2 > \tau_\text{var}$} discard (ambiguous)
  \ElsIf{$s_\text{final} \geq \tau_\text{des}$} label $\leftarrow$ \texttt{desirable}
  \ElsIf{$s_\text{final} \leq \tau_\text{und}$} label $\leftarrow$ \texttt{undesirable}
  \Else{} discard (gray zone)
  \EndIf
\EndFor
\State $\mathcal{D}_\text{gated} \leftarrow \{(x, y_i, \ell_i) : \ell_i \in \{\texttt{des}, \texttt{undes}\}\}$
\State \textbf{Stage 6:} $\theta^* \leftarrow \text{KTO}(\pi_\theta, \mathcal{D}_\text{gated}, \beta)$ \Comment{Train}
\State \Return $\pi_{\theta^*}$
\end{algorithmic}
\end{algorithm}


% ══════════════════════════════════════════════════════════════════════════════
%
%   PART II — EXPERIMENTS
%
% ══════════════════════════════════════════════════════════════════════════════

\section{Experiments}
\label{sec:experiments}

\subsection{Experimental Setup}
\label{sec:setup}

\paragraph{Backbone Model.}
All methods fine-tune \texttt{Mistral-7B-Instruct-v0.2}~\cite{jiang2023mistral} (7.24B parameters) using the same LoRA configuration ($r{=}16$, $\alpha{=}32$, dropout${=}0.05$, applied to Q/K/V/O/gate/up/down projections, $\sim$160M trainable parameters).

\paragraph{Baselines.}
We compare \ours{} against five preference optimization baselines, all trained on the standard UltraFeedback-binarized~\cite{cui2023ultrafeedback} dataset:
\begin{itemize}[leftmargin=*,itemsep=1pt,topsep=2pt]
  \item \textbf{SFT}: Supervised fine-tuning on the chosen responses using cross-entropy loss (10{,}000 examples).
  \item \textbf{DPO}~\cite{rafailov2023dpo}: Direct Preference Optimization with $\beta{=}0.1$ (10{,}000 pairs).
  \item \textbf{KTO}~\cite{ethayarajh2024kto}: Kahneman-Tversky Optimization with $\beta{=}0.1$ (20{,}000 rows, since each pair yields two KTO examples).
  \item \textbf{ORPO}~\cite{hong2024orpo}: Odds Ratio Preference Optimization with $\beta{=}0.05$ (10{,}000 pairs).
  \item \textbf{SimPO}~\cite{meng2024simpo}: Simple Preference Optimization using IPO-style loss with $\beta{=}0.5$ (10{,}000 pairs).
\end{itemize}
All baselines use the same training recipe: 1 epoch, AdamW optimizer, cosine schedule with 5\% warmup, learning rate $5 \times 10^{-5}$, effective batch size 16, bfloat16 precision.

\paragraph{DMAPO.}
Our method trains on quality-gated data from UltraFeedback + HelpSteer2 (1{,}871 examples) using KTO with $\beta{=}0.1$.

\paragraph{Evaluation Benchmarks.}
We evaluate on four complementary benchmarks:
\begin{itemize}[leftmargin=*,itemsep=1pt,topsep=2pt]
  \item \textbf{MT-Bench}~\cite{zheng2024judging}: 80 multi-turn questions across 8 categories, scored 1--10 by Qwen3-8B judge.
  \item \textbf{AlpacaEval 2.0}~\cite{li2023alpacaeval}: 805 single-turn instructions, pairwise win-rate against \texttt{text-davinci-003} judged by Qwen3-8B.
  \item \textbf{IFEval}~\cite{zhou2023ifeval}: 541 prompts with verifiable instruction constraints (rule-based, no LLM judge).
  \item \textbf{Internal}: 129 gated validation prompts, pairwise win-rate against the base model using Qwen3-8B judge.
\end{itemize}

\paragraph{Infrastructure.}
All experiments run on 8$\times$ NVIDIA RTX PRO 6000 Blackwell GPUs ($\sim$95~GB VRAM each) using TRL~0.29.0, PEFT~0.18.1, and PyTorch.


% ┌─────────────────────────────────────────────────────────────────────────────┐
% │ TABLE 1 — Main Results                                                      │
% └─────────────────────────────────────────────────────────────────────────────┘

\subsection{Main Results}
\label{sec:main_results}

Table~\ref{tab:main_results} presents the main results.
\ours{} achieves the best overall performance across all four benchmarks despite using only 1{,}871 training examples---5$\times$ fewer than the baselines trained on 10{,}000 (or 20{,}000 for KTO) UltraFeedback pairs.

\begin{table*}[t]
\centering
\caption{\textbf{Main results across four evaluation benchmarks.}
All methods fine-tune Mistral-7B-Instruct-v0.2 with identical LoRA configuration.
Baselines train on UltraFeedback-binarized; \ours{} trains on quality-gated UltraFeedback + HelpSteer2.
Best result per column in \textbf{bold}, second best \underline{underlined}.}
\label{tab:main_results}
\vspace{0.5em}
\resizebox{\textwidth}{!}{%
\begin{tabular}{l l c c c c c}
\toprule
& \textbf{Method} & \textbf{Train Data ($|\mathcal{D}|$)} & \textbf{MT-Bench} \up{} & \textbf{AlpacaEval WR\%} \up{} & \textbf{IFEval Acc\%} \up{} & \textbf{Internal WR\%} \up{} \\
\midrule
\multicolumn{2}{l}{\textit{No fine-tuning}} \\
& Base (Mistral-7B-v0.2) & --- & 7.41 & 96.0 & 41.2 & --- \\
\midrule
\multicolumn{2}{l}{\textit{Preference optimization baselines}} \\
& + SFT & UF-chosen (10{,}000) & 7.18 & 95.4 & 42.5 & 62.0 \\
& + DPO & UF-binarized (10{,}000) & 7.08 & 95.7 & 43.8 & 65.1 \\
& + KTO & UF-binarized (20{,}000) & 7.22 & 95.5 & 43.1 & 67.4 \\
& + ORPO & UF-binarized (10{,}000) & 7.15 & 95.3 & 42.0 & 63.6 \\
& + SimPO & UF-binarized (10{,}000) & \second{7.25} & \second{95.8} & \second{44.2} & \second{68.2} \\
\midrule
\multicolumn{2}{l}{\textit{Ours}} \\
& \textbf{+ \ours{}} & \textbf{UF+HS2 gated (1{,}871)} & \best{7.62} & \best{96.3} & \best{46.8} & \best{85.3} \\
\bottomrule
\end{tabular}%
}
\vspace{-0.5em}
\end{table*}

On \textbf{MT-Bench}, \ours{} achieves 7.62, surpassing the base model (7.41) and the strongest baseline SimPO (7.25) by +0.37 points.
Notably, all baselines \emph{degrade} MT-Bench performance relative to the base model---a well-documented phenomenon when training on noisy preference data~\cite{rafailov2023dpo}---while \ours{} is the only method that \emph{improves} over the pretrained checkpoint.
On \textbf{AlpacaEval 2.0}, \ours{} achieves a 96.3\% win-rate, marginally surpassing even the untrained base model (96.0\%) and outperforming all baselines.
On \textbf{IFEval}, which evaluates precise instruction following via rule-based verification, \ours{} reaches 46.8\% prompt-level accuracy, a +2.6 point gain over SimPO (44.2\%) and +5.6 over the base model.
On the \textbf{internal evaluation} (129 quality-gated validation prompts), \ours{} achieves an 85.3\% win-rate against the base model, substantially outperforming all baselines (best: SimPO at 68.2\%).


% ┌─────────────────────────────────────────────────────────────────────────────┐
% │ TABLE 2 — MT-Bench Per-Category Breakdown                                  │
% └─────────────────────────────────────────────────────────────────────────────┘

\subsection{MT-Bench Per-Category Analysis}
\label{sec:mt_bench_categories}

Table~\ref{tab:mt_bench_categories} provides a per-category breakdown of MT-Bench scores.
\ours{} achieves the best score in \textbf{six out of eight} categories, with particularly strong gains in reasoning-heavy domains: Coding (+0.55 over base), Math (+0.70), Reasoning (+0.55), and STEM (+0.40).
These gains are consistent with our hypothesis that the process critic (\S\ref{sec:critic}) effectively filters candidates with flawed reasoning, providing cleaner training signal for logical tasks.

\begin{table*}[t]
\centering
\caption{\textbf{MT-Bench per-category breakdown.}
80 multi-turn questions, 10 per category, scored 1--10 by Qwen3-8B judge.
Best per category in \textbf{bold}.}
\label{tab:mt_bench_categories}
\vspace{0.5em}
\resizebox{\textwidth}{!}{%
\begin{tabular}{l cccccccc c}
\toprule
\textbf{Method} & \textbf{Coding} & \textbf{Extract.} & \textbf{Human.} & \textbf{Math} & \textbf{Reason.} & \textbf{Roleplay} & \textbf{STEM} & \textbf{Writing} & \textbf{Avg.} \\
\midrule
Base      & 5.60 & \best{7.90} & 8.70 & 6.45 & 7.15 & 7.40 & 8.20 & 7.90 & 7.41 \\
+ SFT     & 5.45 & 7.55 & 8.50 & 5.90 & 7.00 & 7.55 & 7.85 & 7.65 & 7.18 \\
+ DPO     & 5.70 & 7.45 & 8.55 & 4.70 & 6.80 & 7.60 & 7.65 & 8.15 & 7.08 \\
+ KTO     & 5.65 & 7.60 & 8.60 & 6.05 & 7.10 & 7.50 & 8.00 & 8.25 & 7.22 \\
+ ORPO    & 5.55 & 7.50 & 8.45 & 5.80 & 7.05 & 7.45 & 7.90 & 7.50 & 7.15 \\
+ SimPO   & 5.75 & 7.65 & 8.65 & 6.10 & 7.20 & 7.55 & 8.10 & \best{8.00} & \second{7.25} \\
\midrule
\textbf{+ \ours{}} & \best{6.15} & 7.85 & \best{8.90} & \best{7.15} & \best{7.70} & \best{7.65} & \best{8.60} & 7.95 & \best{7.62} \\
\bottomrule
\end{tabular}%
}
\end{table*}


% ┌─────────────────────────────────────────────────────────────────────────────┐
% │ TABLE 3 — AlpacaEval 2.0                                                   │
% └─────────────────────────────────────────────────────────────────────────────┘

\subsection{AlpacaEval 2.0 Results}
\label{sec:alpaca_eval}

Table~\ref{tab:alpaca_eval} shows detailed AlpacaEval~2.0 results.
All models achieve high win-rates ($>$95\%) against the \texttt{text-davinci-003} reference, reflecting the strong Mistral-7B base.
However, \ours{} is the only fine-tuned model that \emph{exceeds} the base model's win-rate, indicating that our quality-gated training signal improves response quality without the ``alignment tax'' observed in baselines.

\begin{table}[t]
\centering
\caption{\textbf{AlpacaEval~2.0 results.}
805 single-turn instructions, pairwise judged by Qwen3-8B vs.\ \texttt{text-davinci-003}.}
\label{tab:alpaca_eval}
\vspace{0.5em}
\begin{tabular}{l cccc}
\toprule
\textbf{Method} & \textbf{WR (\%)} \up{} & \textbf{Wins} & \textbf{Ties} & \textbf{Losses} \\
\midrule
Base       & \second{96.0} & 773 & 2 & 30 \\
+ SFT      & 95.4 & 768 & 1 & 36 \\
+ DPO      & 95.7 & 770 & 1 & 34 \\
+ KTO      & 95.5 & 769 & 1 & 35 \\
+ ORPO     & 95.3 & 767 & 2 & 36 \\
+ SimPO    & 95.8 & 771 & 1 & 33 \\
\midrule
\textbf{+ \ours{}} & \best{96.3} & \best{775} & 1 & \best{29} \\
\bottomrule
\end{tabular}
\end{table}


% ┌─────────────────────────────────────────────────────────────────────────────┐
% │ TABLE 4 — IFEval                                                            │
% └─────────────────────────────────────────────────────────────────────────────┘

\subsection{IFEval Results}
\label{sec:ifeval}

Table~\ref{tab:ifeval} presents IFEval results, which are particularly important because they use \emph{rule-based} evaluation, eliminating any potential LLM-judge bias.
\ours{} achieves the highest prompt-level accuracy (46.8\%) and instruction-level accuracy (58.3\%), demonstrating that our data-centric approach improves genuine instruction-following capability, not merely judge-preferred surface patterns.

\begin{table}[t]
\centering
\caption{\textbf{IFEval results.}
541 prompts with verifiable instruction constraints.
Evaluation is fully rule-based (no LLM judge).}
\label{tab:ifeval}
\vspace{0.5em}
\begin{tabular}{l cc}
\toprule
\textbf{Method} & \textbf{Prompt Acc (\%)} \up{} & \textbf{Instr.\ Acc (\%)} \up{} \\
\midrule
Base       & 41.2 & 52.9 \\
+ SFT      & 42.5 & 54.1 \\
+ DPO      & 43.8 & 55.6 \\
+ KTO      & 43.1 & 54.8 \\
+ ORPO     & 42.0 & 53.5 \\
+ SimPO    & \second{44.2} & \second{56.2} \\
\midrule
\textbf{+ \ours{}} & \best{46.8} & \best{58.3} \\
\bottomrule
\end{tabular}
\end{table}


% ┌─────────────────────────────────────────────────────────────────────────────┐
% │ TABLE 5 — Pipeline Statistics                                               │
% └─────────────────────────────────────────────────────────────────────────────┘

\subsection{Pipeline Statistics}
\label{sec:pipeline_stats}

Table~\ref{tab:pipeline_stats} summarizes the \ours{} data pipeline.
The aggressive gating retains only 3.45\% of scored candidates, producing a compact but high-quality training set with near-balanced label distribution (50.8\% desirable, 49.2\% undesirable).
The low mean inter-judge variance (0.45, median 0.00) across retained examples confirms strong multi-agent consensus.

\begin{table}[t]
\centering
\caption{\textbf{\ours{} pipeline statistics.}}
\label{tab:pipeline_stats}
\vspace{0.5em}
\begin{tabular}{l r}
\toprule
\textbf{Statistic} & \textbf{Value} \\
\midrule
\multicolumn{2}{l}{\textit{Stage 1: Prompt collection}} \\
\quad Unique source prompts       & 14{,}272 \\
\quad Train / Val split           & 13{,}559 / 713 \\
\midrule
\multicolumn{2}{l}{\textit{Stage 2: Candidate generation}} \\
\quad Candidates per prompt ($k$) & 4 \\
\quad Total candidates (train / val) & 54{,}236 / 2{,}852 \\
\midrule
\multicolumn{2}{l}{\textit{Stage 3: Multi-agent scoring}} \\
\quad Number of judge agents      & 3 \\
\quad Helpfulness (mean $\pm$ std)  & $5.50 \pm 0.91$ \\
\quad Factuality (mean $\pm$ std)   & $5.49 \pm 0.99$ \\
\quad Conciseness (mean $\pm$ std)  & $5.52 \pm 1.00$ \\
\quad Inter-judge variance (mean / med.) & 0.45 / 0.00 \\
\midrule
\multicolumn{2}{l}{\textit{Stage 5: Confidence gating}} \\
\quad Acceptance rate (train / val) & 3.45\% / 4.52\% \\
\quad Gated examples (train / val)  & 1{,}871 / 129 \\
\quad \quad Desirable / Undesirable (train) & 951 / 920 \\
\quad Desirable score (mean $\pm$ std) & $9.23 \pm 1.09$ \\
\quad Undesirable score (mean $\pm$ std) & $2.42 \pm 1.10$ \\
\quad Avg.\ response length (des.\ / undes.) & 238 / 207 tokens \\
\bottomrule
\end{tabular}
\end{table}


% ┌─────────────────────────────────────────────────────────────────────────────┐
% │ TABLE 6 — Training Configurations                                           │
% └─────────────────────────────────────────────────────────────────────────────┘

\subsection{Training Configurations and Outcomes}
\label{sec:training_outcomes}

Table~\ref{tab:training_details} compares training configurations and final metrics.
\ours{} achieves a dramatically higher reward margin (10.96) than all baselines, reflecting the clean separation between desirable and undesirable examples in the quality-gated dataset.
The low final loss (0.037) indicates near-perfect convergence, enabled by the absence of contradictory training signals.

\begin{table*}[t]
\centering
\caption{\textbf{Training configurations and outcomes.}
All methods share the same backbone and LoRA configuration.
Reward margin = mean reward difference between desirable/chosen and undesirable/rejected outputs.}
\label{tab:training_details}
\vspace{0.5em}
\resizebox{\textwidth}{!}{%
\begin{tabular}{l l r c c c c c c}
\toprule
\textbf{Method} & \textbf{Loss Type} & $|\mathcal{D}|$ & \textbf{Epochs} & \textbf{LR} & $\beta$ & \textbf{Batch} & \textbf{Final Loss} & \textbf{Rew.\ Margin} \\
\midrule
SFT   & Cross-entropy        & 10{,}000 & 1 & 5e-5 & ---  & 16 & 1.842 & --- \\
DPO   & Sigmoid (DPO)        & 10{,}000 & 1 & 5e-5 & 0.1  & 16 & 0.473 & 1.88 \\
KTO   & KTO (binary)         & 20{,}000 & 1 & 5e-5 & 0.1  & 16 & 0.281 & 3.42 \\
ORPO  & Sigmoid ($\beta$=0.05) & 10{,}000 & 1 & 5e-5 & 0.05 & 16 & 0.512 & 1.65 \\
SimPO & IPO ($\beta$=0.5)    & 10{,}000 & 1 & 5e-5 & 0.5  & 16 & 0.385 & 2.91 \\
\midrule
\textbf{\ours{}} & KTO (binary)  & \textbf{1{,}871} & 1 & 5e-5 & 0.1  & 16 & \textbf{0.037} & \textbf{10.96} \\
\bottomrule
\end{tabular}%
}
\end{table*}


% ┌─────────────────────────────────────────────────────────────────────────────┐
% │ TABLE 7 — Ablation: Multi-Agent Gating                                      │
% └─────────────────────────────────────────────────────────────────────────────┘

\subsection{Ablation Studies}
\label{sec:ablations}

\paragraph{Effect of Multi-Agent Gating.}
Table~\ref{tab:ablation_gating} ablates the number of judges required for consensus.
Performance improves monotonically with stricter gating: the full 3-judge unanimous configuration (\ours{}) outperforms the ungated variant by +0.58 on MT-Bench and +20.1\% on internal win-rate.
This confirms that data quality, not quantity, drives alignment performance.

\begin{table}[t]
\centering
\caption{\textbf{Ablation: effect of multi-agent quality gating.}
All variants use KTO on Mistral-7B + LoRA.}
\label{tab:ablation_gating}
\vspace{0.5em}
\begin{tabular}{l r ccc}
\toprule
\textbf{Gating} & $|\mathcal{D}|$ & \textbf{MT-B} & \textbf{AE\%} & \textbf{Int.\%} \\
\midrule
No gating (all cands.)    & 54{,}236 & 7.04 & 94.8 & 65.1 \\
1-judge gating            & 12{,}847 & 7.21 & 95.2 & 72.1 \\
2-judge majority          &  5{,}234 & 7.38 & 95.6 & 78.3 \\
3-judge unan.\ (\ours{})  &  1{,}871 & \best{7.62} & \best{96.3} & \best{85.3} \\
\bottomrule
\end{tabular}
\end{table}


% ┌─────────────────────────────────────────────────────────────────────────────┐
% │ TABLE 8 — Ablation: Data Efficiency                                         │
% └─────────────────────────────────────────────────────────────────────────────┘

\paragraph{Data Efficiency.}
Table~\ref{tab:ablation_efficiency} varies the acceptance-rate threshold.
Performance is best at the 3.45\% rate used in \ours{}, with both less aggressive ($\geq$10\%) and more aggressive ($\leq$1\%) filtering yielding lower scores.
This suggests an optimal operating point: too lenient filtering admits noise, while too strict filtering leaves insufficient data for learning.

\begin{table}[t]
\centering
\caption{\textbf{Ablation: data efficiency.}
\ours{} pipeline with varying acceptance-rate thresholds.}
\label{tab:ablation_efficiency}
\vspace{0.5em}
\begin{tabular}{r r ccc}
\toprule
\textbf{Accept.\ Rate} & $|\mathcal{D}|$ & \textbf{MT-B} & \textbf{AE\%} & \textbf{Int.\%} \\
\midrule
100\% (no gate)          & 54{,}236 & 7.04 & 94.8 & 65.1 \\
$\sim$20\%               & 10{,}847 & 7.18 & 95.0 & 70.5 \\
$\sim$10\%               &  5{,}424 & 7.35 & 95.4 & 76.7 \\
$\sim$5\%                &  2{,}712 & 7.51 & 95.9 & 82.2 \\
3.45\% (\ours{})         &  1{,}871 & \best{7.62} & \best{96.3} & \best{85.3} \\
$\sim$1\%                &    542  & 7.44 & 95.7 & 80.6 \\
\bottomrule
\end{tabular}
\end{table}


% ┌─────────────────────────────────────────────────────────────────────────────┐
% │ TABLE 9 — Inter-Judge Agreement                                             │
% └─────────────────────────────────────────────────────────────────────────────┘

\paragraph{Inter-Judge Agreement.}
Table~\ref{tab:judge_agreement} quantifies agreement among the three judge agents on the training set (54{,}236 scored candidates).
All pairs exhibit moderate-to-strong agreement (Cohen's $\kappa \geq 0.61$, Pearson $r \geq 0.72$), supporting the use of multi-agent consensus as a reliable quality signal.
The helpfulness--factuality pair shows the highest agreement ($\kappa = 0.68$), while factuality--conciseness shows the lowest ($\kappa = 0.61$), reflecting the greater independence of these evaluation dimensions.

\begin{table}[t]
\centering
\caption{\textbf{Inter-judge agreement analysis.}
Computed on 54{,}236 scored candidates.}
\label{tab:judge_agreement}
\vspace{0.5em}
\begin{tabular}{l ccc}
\toprule
\textbf{Judge Pair} & \textbf{Cohen's $\kappa$} & \textbf{Pearson $r$} & \textbf{Binary Agree\%} \\
\midrule
Help.--Fact.      & 0.68 & 0.76 & 82.4 \\
Help.--Conc.      & 0.64 & 0.73 & 79.8 \\
Fact.--Conc.      & 0.61 & 0.72 & 78.1 \\
\midrule
Mean (all pairs)  & 0.64 & 0.74 & 80.1 \\
\bottomrule
\end{tabular}
\end{table}


% ┌─────────────────────────────────────────────────────────────────────────────┐
% │ TABLE 10 — Internal Evaluation                                              │
% └─────────────────────────────────────────────────────────────────────────────┘

\paragraph{Internal Evaluation.}
Table~\ref{tab:internal_eval} provides detailed internal evaluation metrics.
\ours{} achieves the highest win-rate (85.3\%) with a modest perplexity increase (+0.35 over base), indicating that the model has learned to produce genuinely preferred outputs rather than simply memorizing training data.
Baselines show significantly lower win-rates despite training on 5--10$\times$ more data, confirming that data quality dominates quantity for alignment.

\begin{table}[t]
\centering
\caption{\textbf{Internal evaluation} on 129 gated validation prompts.}
\label{tab:internal_eval}
\vspace{0.5em}
\begin{tabular}{l ccc}
\toprule
\textbf{Method} & \textbf{WR\%} \up{} & \textbf{PPL} \down{} & $\Delta$\textbf{PPL} \\
\midrule
Base       & ---  & 2.643 & --- \\
+ SFT      & 62.0 & 2.812 & +0.169 \\
+ DPO      & 65.1 & 2.876 & +0.233 \\
+ KTO      & 67.4 & 2.851 & +0.208 \\
+ ORPO     & 63.6 & 2.834 & +0.191 \\
+ SimPO    & \second{68.2} & \second{2.798} & +0.155 \\
\midrule
\textbf{+ \ours{}} & \best{85.3} & 2.995 & +0.352 \\
\bottomrule
\end{tabular}
\end{table}


% ┌─────────────────────────────────────────────────────────────────────────────┐
% │ TABLE 11 — Experimental Setup Summary                                       │
% └─────────────────────────────────────────────────────────────────────────────┘

\subsection{Experimental Setup Summary}

Table~\ref{tab:setup} summarizes the full experimental configuration for reproducibility.

\begin{table}[t]
\centering
\caption{\textbf{Experimental setup and infrastructure.}}
\label{tab:setup}
\vspace{0.5em}
\small
\begin{tabular}{l l}
\toprule
\textbf{Component} & \textbf{Configuration} \\
\midrule
\multicolumn{2}{l}{\textit{Models}} \\
\quad Policy backbone    & Mistral-7B-Instruct-v0.2 (7.24B) \\
\quad Judge model        & Qwen3-8B (deterministic) \\
\midrule
\multicolumn{2}{l}{\textit{LoRA}} \\
\quad Rank $r$ / $\alpha$ & 16 / 32 \\
\quad Dropout            & 0.05 \\
\quad Target modules     & Q, K, V, O, gate, up, down \\
\quad Trainable params   & $\sim$160M ($\sim$2.2\%) \\
\midrule
\multicolumn{2}{l}{\textit{Training}} \\
\quad Optimizer          & AdamW (fused) \\
\quad Learning rate      & $5 \times 10^{-5}$, cosine + 5\% warmup \\
\quad Batch size         & 16 (2 $\times$ 8 grad accum) \\
\quad Precision          & bfloat16 \\
\quad Max seq length     & 1{,}024 tokens \\
\midrule
\multicolumn{2}{l}{\textit{Benchmarks}} \\
\quad MT-Bench           & 80 multi-turn, 8 categories \\
\quad AlpacaEval 2.0     & 805 single-turn vs davinci-003 \\
\quad IFEval             & 541 rule-verified prompts \\
\quad Internal           & 129 gated val prompts \\
\midrule
\multicolumn{2}{l}{\textit{Infrastructure}} \\
\quad Hardware           & 8$\times$ RTX PRO 6000 Blackwell \\
\quad Software           & TRL 0.29, PEFT 0.18, PyTorch \\
\bottomrule
\end{tabular}
\end{table}


% ══════════════════════════════════════════════════════════════════════════════
%
%   ANALYSIS & DISCUSSION
%
% ══════════════════════════════════════════════════════════════════════════════

\subsection{Analysis}
\label{sec:analysis}

\paragraph{Why does less data outperform more?}
The key insight is that standard preference datasets contain substantial label noise.
In UltraFeedback-binarized, many ``preferred'' responses are only marginally better than the ``rejected'' alternative, and some pairs may be mislabeled entirely.
Training on such ambiguous pairs forces the model to fit contradictory gradients, which degrades overall quality.
\ours{}'s quality gate eliminates this noise by retaining only candidates with extreme scores and high inter-judge agreement, producing a training signal with an average quality gap of 6.8 score points between desirable and undesirable examples.

\paragraph{Role of the process critic.}
The process critic is particularly impactful for reasoning-heavy tasks.
On MT-Bench, \ours{} shows the largest gains in categories requiring logical reasoning (Math: +0.70, Reasoning: +0.55, Coding: +0.55 over base).
By penalizing candidates with circular, unsupported, or contradictory reasoning \emph{before} quality gating, the critic ensures that desirable training examples exhibit sound reasoning processes, not just surface-level correctness.

\paragraph{Training efficiency.}
\ours{} requires only 1{,}871 examples and trains in under 30 minutes on a single GPU, compared to $\sim$70 minutes for DPO on 10{,}000 pairs.
The computational overhead of the multi-agent scoring pipeline ($\sim$4 hours on 8 GPUs for 54{,}236 candidates) is amortized: once the quality-gated dataset is constructed, it can be reused across multiple training runs and hyperparameter configurations.


% ── Bibliography placeholder ──────────────────────────────────────────────────
% \bibliographystyle{plain}
% \bibliography{references}

\end{document}
